\chapter{关系认知}

\section{关系状态与关于某人的记忆检索}

关系领域（Relationship）维护主体与特定个体或群体之间的长期、关系特定认知。关于某人的记忆检索为关系判断提供经历材料；关系状态组织双方角色、互动期待、共享经历引用、相处边界和关系特定解释。\cite{nous-research-package,baldwin1992}

关系领域的最小对象和字段范围，应由身份、修订、证据、生命周期、查询和写入合同中的稳定差异决定。

\section{证据与关系更新}

互动次数作为经历证据之一，由关系领域结合来源和语境解释。身份绑定依据跨平台证据；显示名称作为来源字段保留。关系变化先形成带来源信息的提案，再由关系领域核验。修订或撤回关系判断时，底层消息事实和观察记录保持原样。

跨平台身份关联应有证据支持。来源中出现别名、昵称或账户标识时，应保存对应区域和身份绑定历史；关系语义与外部授权分别归属适当的领域。

\section{当前状态}

\StatusTag{目标}{关系认知是完整主体档案中的独立高级领域。}

\StatusTag{待定}{关系对象、状态集合和写入范围需通过后续研究确定。}

\StatusTag{待实现}{当前检视的 Heptalogos 主体包和 Nous Wave 源树尚无关系认知领域；其语义仍处于目标设计阶段。}
